% homework/hw09x/hw09x.tex
% Year 7 - Extra Homework: where the point goes.
%
% Goes out alongside Homework 9. Same framing rule as hw03x through hw08x:
% nothing on the sheet says easy, basic, extra help, catch-up or revision, and
% no question is marked as being for anybody in particular. The cover states a
% fact, and it is a true one:
%
%   There is nothing new to calculate in Unit 4. 3.2 x 7 is 32 x 7. The one
%   new question is WHERE THE POINT GOES.
%
% THAT IS THE WHOLE DESIGN. Every method in Unit 4 is printed here as a frame
% the student fills in: column grids with the point column already drawn and
% the zeros to be written tinted orange; a point-out / point-back table whose
% scaffold fades row by row; division frames with the point already placed in
% the answer. Filling the boxes IS the method, so it cannot be skipped.
%
% hw08x was place value and the point never moving. This sheet does not repeat
% it: there is no place-value table and no x 10 / x 100 here.
%
% QUESTION NUMBERS ARE X1..X19. Homework 9 has Sections A-F and the two
% sheets go home in the same envelope.
%
% FRENCH COPY. \fr{...} prints a grey French line under the English and
% \frin{...} a short French gloss in brackets. Both are EMPTY in the English
% copy, so the English layout is the one that was checked to eight pages. The
% French copy is built with \FRENCH defined (see README) and runs longer. The
% English key words stay English inside the French lines, in blue, because the
% English is still what is being learned. Numbers keep the English decimal
% point, as the class does.
%
% House style is shared: see homework/hw-style.tex. Method: the new-homework
% skill. Source is pure ASCII on purpose: French accents are written \'e etc.
%
% Build:  pwsh .claude/skills/new-homework/scripts/build-hw.ps1 hw09x -French

\documentclass[12pt,a4paper]{article}

% -- Answer key switch -------------------------------------------------------
\newif\ifkey
\ifdefined\TEACHER\keytrue\else
  \keyfalse        % \keyfalse = student copy   |   \keytrue = teacher copy
\fi

% -- Language switch ---------------------------------------------------------
\newif\iffr
\ifdefined\FRENCH\frtrue\fi

% -- House style -------------------------------------------------------------
% homework/hw-style.tex
% Shared house style for every Year 7 homework packet. \input this from the
% packet file, right after \documentclass and the answer-key switch.
%
% Do not fork this file per packet. If a packet needs a different look, the
% question is whether every packet should have it.
%
% The choices here are deliberate and were arrived at by printing and looking:
%   * No solid dark banners. Every header is a light tint with a coloured spine
%     on the left, so colour reads clearly without flooding a school printer.
%   * No marks and no mark totals. These are practice, not exam papers.
%   * Lato at 12pt with open leading, plus mathastext so inline maths is Lato
%     too. Without mathastext every "-6 + 4" is Computer Modern serif and the
%     sheet looks half-typeset.
%   * Writing lines are spaced for Year 7 handwriting, not adult margin notes.
%   * Pages fill. \needspace is redefined below so a heading no longer drags
%     the page break up to itself, and every heading is glued to the box under
%     it. Packets built before 2026-09-30 were laid out under the old rules and
%     will reflow if rebuilt -- re-read every page if you rebuild one.
%
% Provides: \wlines \blank \tickbox \sep \qmk-free marks (none) \hwsection
% \qhead \expr, the workspace box, and the remember / wordhelp / activity /
% yourturn coloured boxes. Palette matches the lesson decks exactly.

% -- Packages ----------------------------------------------------------------
\usepackage[T1]{fontenc}
\usepackage[utf8]{inputenc}
\usepackage[default]{lato}
\usepackage[a4paper,top=1.7cm,bottom=1.7cm,left=2.0cm,right=2.0cm,headheight=15pt]{geometry}
\usepackage{amsmath,amssymb}
\usepackage{xcolor}
\usepackage[most]{tcolorbox}
\usepackage{enumitem}
\usepackage{tikz}
\usetikzlibrary{arrows.meta}
\usepackage{array}
\usepackage{tabularx}
\usepackage{fancyhdr}
\usepackage{multicol}
\usepackage{needspace}
\usepackage{graphicx}

% Make maths use Lato too. Without this, every "-6 + 4" on the page is set in
% Computer Modern serif and the sheet looks half-typeset. Must come last.
\usepackage[italic]{mathastext}

% -- \needspace, without the hole it used to leave ----------------------------
% needspace.sty writes "\vskip 0pt plus X \penalty-100" in front of every
% heading. A page that ends anywhere else is underfull with no stretch, so TeX
% scores it b=10000; a page that ends AT a heading gets that X of stretch and a
% finite badness, and wins. So any heading in roughly the bottom half of a page
% claimed the break, and the question under it moved over whole, leaving a hole
% the size of the question. \tracingpages showed it exactly (b=4531 at the
% heading, b=10000 at every later break) while building homework/q1rev, and it
% is where the half-empty pages in HW 1-8 came from.
% This keeps only the \penalty9999, which still refuses to put a heading in the
% last X of a page. Keeping a heading WITH the box under it is \nobreak's job
% now -- see \hwsection and \qhead below.
\makeatletter
\renewcommand{\needspace}[1]{\begingroup\setlength{\dimen@}{#1}%
  \vskip\dimen@\penalty9999\vskip-\dimen@\endgroup}
\makeatother

\linespread{1.06}\selectfont

% -- The Learner's Book palette (same hex values as the lesson decks) ---------
\definecolor{cbteal}{HTML}{0087A8}
\definecolor{cbpurple}{HTML}{5C2483}
\definecolor{cborange}{HTML}{C25E12}
\definecolor{cbgreen}{HTML}{4A8B23}
\definecolor{cbred}{HTML}{C8102E}
\definecolor{cbblue}{HTML}{1A5FA8}
\definecolor{rulegray}{HTML}{9AA5AD}

% -- Answer lines and blanks -------------------------------------------------
% \wlines{n} draws n ruled writing lines, spaced wide enough for Year 7
% handwriting rather than for an adult with a fine pen.
\makeatletter
\newcommand{\wlines}[1]{%
  \par\vspace{0.75em}%
  \@tempcnta=0\relax
  \@whilenum\@tempcnta<#1\do{%
    \hbox to \linewidth{\color{rulegray}\leaders\hrule height 0.5pt\hfill}%
    \vspace{1.5em}%
    \advance\@tempcnta by 1\relax}%
  \par\vspace{-0.8em}%
}
\makeatother

% \blank[width] is an inline gap to write a short answer in.
\newcommand{\blank}[1][2.6cm]{\underline{\hspace{#1}}}

% A boxed space for showing working.
\newtcolorbox{workspace}[1][3.4cm]{%
  colback=white, colframe=rulegray, boxrule=0.5pt, arc=5pt,
  height=#1, left=6pt, right=6pt, top=4pt, bottom=4pt,
  before skip=9pt, after skip=11pt}

% An empty tick box.
\newcommand{\tickbox}{\fbox{\rule{0pt}{1.15em}\rule{1.15em}{0pt}}}

% A middot separator.
\newcommand{\sep}{\,\textperiodcentered\,}

% -- Section and question headings -------------------------------------------
% Light tint plus a fat coloured spine on the left. No solid dark banner: the
% colour reads clearly but barely costs the printer anything.
% needspace keeps a heading out of the last few lines of a page. The \nobreak
% after it keeps it with whatever comes next: a remember or word-help box that
% does not fit takes its heading over with it, instead of leaving the heading
% alone at the foot of the page. (The old needspace did that by accident, and
% left a hole every time; see above.)
\newcommand{\hwsection}[2][cbteal]{%
  \needspace{8\baselineskip}%
  \par\vspace{13pt}%
  \noindent
  \begin{tcolorbox}[enhanced, nobeforeafter, width=\textwidth,
      colback=#1!7, colframe=#1!7, arc=5pt,
      borderline west={5pt}{0pt}{#1},
      left=14pt, right=10pt, top=7pt, bottom=7pt]
    {\large\bfseries\color{#1}#2}
  \end{tcolorbox}%
  \par\nobreak\vspace{9pt}}

% \qhead[colour]{A2}{Moving along the line} -- a question heading.
\newcommand{\qhead}[3][cbteal]{%
  \needspace{6\baselineskip}%
  \par\vspace{11pt}%
  \noindent{\large\bfseries\color{#1}#2}\quad{\large\bfseries #3}\par\nobreak\vspace{4pt}}

% -- Coloured boxes ----------------------------------------------------------
% Shared look: rounded, light fill, coloured rule, coloured title text on a
% pale title strip rather than a dark one.
% NOT breakable. A breakable box that cannot fit its first lines forces its own
% page break, and \nobreak cannot stop a forced break -- so the heading above
% it was left alone at the foot of the page (tested: it happens with \nobreak
% in place). Unbreakable, the page builder moves heading and box together.
% These boxes are a few lines each; none needs to split.
\tcbset{hwbox/.style={
  enhanced, arc=6pt, boxrule=1pt,
  left=11pt, right=11pt, top=8pt, bottom=8pt,
  fonttitle=\bfseries, attach boxed title to top left={xshift=12pt, yshift=-8pt},
  boxed title style={arc=4pt, boxrule=1pt, left=7pt, right=7pt, top=2pt, bottom=2pt},
  top=13pt}}

% Orange: things to remember. Matches the deck's "Write This Down".
\newtcolorbox{remember}[1][Remember from the lesson]{hwbox,
  colback=cborange!6, colframe=cborange, coltitle=cborange,
  colbacktitle=white, title={#1}}

% Blue: English help. The barrier in this class is the wording, not the maths.
\newtcolorbox{wordhelp}[1][Word help]{hwbox,
  colback=cbblue!6, colframe=cbblue, coltitle=cbblue,
  colbacktitle=white, title={#1}}

% Purple: an activity or instruction.
\newtcolorbox{activity}[1][How to do this homework]{hwbox,
  colback=cbpurple!6, colframe=cbpurple, coltitle=cbpurple,
  colbacktitle=white, title={#1}}

% Green: your turn / a friendly nudge.
\newtcolorbox{yourturn}[1][Your turn]{hwbox,
  colback=cbgreen!7, colframe=cbgreen, coltitle=cbgreen!75!black,
  colbacktitle=white, title={#1}}

% -- Shorthands --------------------------------------------------------------
\newcommand{\dC}{\,^{\circ}\mathrm{C}}          % use inside maths: $-8\dC$
\newcommand{\key}[1]{{\color{cborange}\textbf{#1}}}
\newcommand{\eng}[1]{{\color{cbblue}\textbf{#1}}}

\setlist[enumerate]{itemsep=8pt, topsep=5pt, leftmargin=*}
\setlist[itemize]{itemsep=4pt, topsep=4pt, leftmargin=*}
\renewcommand{\arraystretch}{1.45}

% -- An expression the student is asked to work from -------------------------
\newcommand{\expr}[1]{%
  \tcbox[on line, colback=cbgreen!10, colframe=cbgreen!55, boxrule=1pt, arc=4pt,
         left=9pt, right=9pt, top=4pt, bottom=4pt]{\large\bfseries $#1$}}

% -- Running head and foot ---------------------------------------------------
% The packet overrides \fancyhead[L] with its own title.
\pagestyle{fancy}
\fancyhf{}
\fancyhead[L]{\small\color{black!55}Year 7\sep Homework}
\fancyhead[R]{\small\color{black!55}Name: \blank[3.4cm]}
\fancyfoot[C]{\small\color{black!55}Page \thepage}
\renewcommand{\headrulewidth}{0pt}
\renewcommand{\headrule}{\vspace{2pt}\hbox to\headwidth{\color{cbteal!45}\leaders\hrule height 1.2pt\hfill}}



% -- French support lines ----------------------------------------------------
% \fr{...}  a French line of its own, grey and italic, under the English.
% \frin{...} a French gloss in brackets, inline -- for table cells and labels.
% \fq{...}  French quotation marks.
% Both start with \unskip: the line break before \fr is a space, and in the
% English copy it would otherwise sit at the end of the line and move a break.
\newcommand{\fr}[1]{\unskip\iffr\par\vspace{1pt}\noindent{\small\itshape\color{black!60}#1}\par\fi}
\newcommand{\frin}[1]{\unskip\iffr\ {\small\itshape\color{black!60}(#1)}\fi}
\newcommand{\fq}[1]{\guillemotleft\,#1\,\guillemotright}

% -- Per-packet running head -------------------------------------------------
\iffr
  \fancyhead[L]{\small\color{black!55}Year 7\sep Extra Homework\sep English /
    Fran\c{c}ais}
\else
  \fancyhead[L]{\small\color{black!55}Year 7\sep Extra Homework}
\fi

% -- A group of writing lines that will not be split -------------------------
\newcommand{\qlines}[1]{\par\needspace{\numexpr#1+1\relax\baselineskip}\wlines{#1}}

% -- A box to write one digit in ---------------------------------------------
\newcommand{\nbox}{\ \fbox{\rule{0pt}{1.05em}\rule{1.0em}{0pt}}\ }
% The same box, tinted: write a 0 here.
\newcommand{\zbox}{\ \fcolorbox{cborange}{cborange!14}{\rule{0pt}{1.05em}\rule{1.0em}{0pt}}\ }
% A box with no space round it, to sit against a number: 0.6\zb
\newcommand{\zb}{\fcolorbox{cborange}{cborange!14}{\rule{0pt}{0.95em}\rule{0.75em}{0pt}}}

% -- An answer blank on a line of its own, right-aligned ---------------------
\newcommand{\ansline}[1][4.4cm]{\par\vspace{2pt}\hspace*{\fill}\blank[#1]\par}

% -- A ragged-right X column and a centred fixed one -------------------------
\newcolumntype{L}{>{\raggedright\arraybackslash}X}
\newcolumntype{C}[1]{>{\centering\arraybackslash}p{#1}}

% -- A table row tall enough to write in -------------------------------------
\newcommand{\tall}{\rule[-0.6em]{0pt}{1.9em}}

% -- More room in the French copy --------------------------------------------
% Same scheme as hw09.tex: looser leading, taller table rows and working boxes
% a third taller in the French copy only, which also lands it on an even number
% of pages (10). The English copy is untouched.
\iffr
  \linespread{1.10}\selectfont
  \renewcommand{\tall}{\rule[-0.8em]{0pt}{2.4em}}
  \renewtcolorbox{workspace}[1][3.4cm]{%
    colback=white, colframe=rulegray, boxrule=0.5pt, arc=5pt,
    height={1.33\dimexpr #1\relax}, left=6pt, right=6pt, top=4pt, bottom=4pt,
    before skip=9pt, after skip=11pt}
\fi

% ---------------------------------------------------------------------------
% COLUMN GRID. \colgrid{op}{row 1}{row 2}{answer row}
% Each row is four entries: tens, ones, tenths, hundredths. The point column
% sits between ones and tenths and is printed on all three rows, so it cannot
% move. An entry is a digit, nothing, z (an orange box: write a 0 here) or Z
% (an orange box with the 0 already written -- for the worked example).
% ---------------------------------------------------------------------------
\def\zz{z}
\def\ZZ{Z}
\def\pt{.}
\newcommand{\gcell}[4][black]{%
  \edef\gcE{#4}%
  \ifx\gcE\zz
    \filldraw[fill=cborange!14, draw=cborange!80, line width=0.8pt]
      (#2,#3) rectangle ++(0.72,0.78);
  \else\ifx\gcE\ZZ
    \filldraw[fill=cborange!14, draw=cborange!80, line width=0.8pt]
      (#2,#3) rectangle ++(0.72,0.78);
    \node[font=\large\bfseries, text=cborange] at (#2+0.36,#3+0.39) {0};
  \else
    \draw[rulegray, line width=0.6pt] (#2,#3) rectangle ++(0.72,0.78);
    \node[font=\large, text=#1] at (#2+0.36,#3+0.39) {\gcE};
  \fi\fi}
\newcommand{\colgrid}[4]{%
  \begin{tikzpicture}[x=1cm, y=1cm, baseline=(current bounding box.north)]
    \foreach \d [count=\i] in {#2}{%
      \pgfmathsetmacro\xx{(\i-1)*0.72+(\i>2)*0.30}\gcell{\xx}{1.75}{\d}}
    \foreach \d [count=\i] in {#3}{%
      \pgfmathsetmacro\xx{(\i-1)*0.72+(\i>2)*0.30}\gcell{\xx}{0.85}{\d}}
    \foreach \d [count=\i] in {#4}{%
      \pgfmathsetmacro\xx{(\i-1)*0.72+(\i>2)*0.30}\gcell[cbteal]{\xx}{-0.25}{\d}}
    \foreach \y in {1.75,0.85,-0.25}{\fill (1.59,\y+0.13) circle (0.06);}
    \node[font=\Large] at (-0.42,1.24) {$#1$};
    \draw[line width=1.2pt] (-0.75,0.72) -- (3.18,0.72);
  \end{tikzpicture}}

% ---------------------------------------------------------------------------
% DIVISION FRAME. \bus{divisor}{dividend}
% The dividend is a list of entries: a digit, . (the point) or z (an orange
% box: write a 0 here and keep dividing). Above every digit is an empty box for
% the answer; above the point the point is already printed. That is the whole
% of "the point goes straight up", drawn.
% ---------------------------------------------------------------------------
% \bus[np]{..}{..} leaves the point OUT of the answer row: the student puts it
% there. That is X12, the step after X11.
\def\np{np}
\newcommand{\bus}[3][]{%
  \edef\busopt{#1}%
  \begin{tikzpicture}[x=1cm, y=1cm, baseline=(current bounding box.north)]
    \draw[line width=1.2pt] (-0.20,0.86) -- (0,0.86);
    \foreach \d [count=\i] in {#3}{%
      \pgfmathsetmacro\xx{(\i-1)*0.66}%
      \edef\gcE{\d}%
      \draw[line width=1.2pt] (\xx,0.86) -- (\xx+0.66,0.86);
      \ifx\gcE\pt
        \fill (\xx+0.33,0.12) circle (0.06);
        \ifx\busopt\np\else\fill (\xx+0.33,1.10) circle (0.06);\fi
      \else
        \draw[rulegray, line width=0.6pt] (\xx+0.04,0.98) rectangle ++(0.58,0.74);
        \ifx\gcE\zz
          \filldraw[fill=cborange!14, draw=cborange!80, line width=0.8pt]
            (\xx+0.04,0.04) rectangle ++(0.58,0.74);
        \else
          \node[font=\large] at (\xx+0.33,0.41) {\gcE};
        \fi
      \fi}
    \draw[line width=1.2pt] (-0.20,-0.02) .. controls (-0.06,0.42) .. (-0.20,0.86);
    \node[font=\large] at (-0.55,0.41) {#2};
  \end{tikzpicture}}

\begin{document}

% ===========================================================================
% COVER BLOCK
% ===========================================================================
\thispagestyle{fancy}

\begin{tcolorbox}[enhanced, colback=cbgreen!8, colframe=cbgreen!8, arc=7pt,
                  borderline west={6pt}{0pt}{cbgreen},
                  left=18pt, right=14pt, top=11pt, bottom=11pt]
  {\color{cbgreen!75!black}\bfseries YEAR 7\sep UNIT 4\sep EXTRA HOMEWORK}\\[6pt]
  {\color{cbgreen!70!black}\Huge\bfseries Where the Point Goes}\\[10pt]
  {\color{black!70}Comparing\sep adding and subtracting\sep multiplying and
  dividing\sep the words that choose}
  \iffr\\[4pt]{\small\itshape\color{black!60}O\`u placer la virgule~? Comparer,
  additionner, soustraire, multiplier, diviser --- et les mots qui indiquent
  l'op\'eration}\fi
\end{tcolorbox}

\vspace{6pt}
\noindent
\begin{tabularx}{\textwidth}{@{}X X@{}}
  \textbf{Name:} \blank[5.6cm] & \textbf{Class:} \blank[5.6cm] \\
  \textbf{Date given:} \blank[4.6cm] &
  \textbf{Date due:} {\color{cbred}\bfseries Friday 23 October 2026} \\
\end{tabularx}

\vspace{4pt}
\begin{activity}[Why this sheet exists]
  There is \key{nothing new to calculate} in Unit 4. $3.2 \times 7$ is $32 \times 7$,
  and you already know how to do that. $8.24 \div 4$ is short division. The one new
  question is \key{where the point goes}.
  \fr{Dans l'unit\'e 4, il n'y a \textbf{rien de nouveau \`a calculer}. $3.2 \times
  7$, c'est $32 \times 7$, et tu sais d\'ej\`a le faire. La seule nouvelle question,
  c'est~: \textbf{o\`u placer la virgule} (\eng{where the point goes})~? En anglais,
  la virgule s'appelle \eng{the point} et s'\'ecrit avec un point~: 3.25.}

  \medskip
  Every question on this sheet prints its steps as \key{boxes}. Fill the boxes ---
  that \emph{is} the method. An \key{orange box} means: write a \key{0} here.
  \fr{Chaque question pr\'esente ses \'etapes dans des cases~: remplir les cases,
  c'est appliquer la m\'ethode. Une case orange veut dire~: \'ecris un 0 ici.}

  \medskip
  \eng{No calculator.} \frin{Calculatrice interdite.}
\end{activity}

% ===========================================================================
% PART 1 -- WHICH IS BIGGER?
% ===========================================================================
\hwsection[cbgreen]{Part 1 \textperiodcentered\ Which is bigger?\frin{Lequel est le plus grand~?}}

% -- X1 ---------------------------------------------------------------------
\qhead[cbgreen]{X1}{Two parts of a number}

\noindent Write the \textbf{whole-number part} and the \textbf{decimal part}. The first
one is done for you.
\fr{\'Ecris la partie enti\`ere (\eng{whole-number part}) et la partie d\'ecimale
(\eng{decimal part}) de chaque nombre. La premi\`ere ligne est faite pour toi.}

\vspace{4pt}
\begin{center}
\begin{tabular}{|C{3.0cm}|C{4.4cm}|C{4.4cm}|}
  \hline
  \textbf{Number} & \textbf{Whole-number part} & \textbf{Decimal part} \\ \hline
  $3.25$ \tall  & 3 & .25 \\ \hline
  $14.6$ \tall  & & \\ \hline
  $0.38$ \tall  & & \\ \hline
  $7.05$ \tall  & & \\ \hline
  $120.9$ \tall & & \\ \hline
\end{tabular}
\end{center}

% -- X2 ---------------------------------------------------------------------
% The deck's rule as a frame. The orange boxes are the zeros, and they are
% printed exactly where they go, so "fill with zeros" is a thing the hand does.
\qhead[cbgreen]{X2}{Fill with zeros, then compare}

\begin{remember}[Compare decimals]
  First look at the \key{whole-number parts}. If they are the same, \key{fill with
  zeros} so both numbers have the same number of digits after the point. Then
  compare. \quad $3.3 = 3.30$, \; and $3.30 > 3.25$.
  \fr{Compare d'abord les parties enti\`eres. Si elles sont \'egales, ajoute des
  z\'eros (\eng{fill with zeros}) pour avoir autant de chiffres apr\`es la virgule,
  puis compare.}
\end{remember}

\vspace{2pt}
\noindent Write a 0 in each orange box. Then write the \textbf{bigger} number. The
first one is done for you.
\fr{\'Ecris un 0 dans chaque case orange, puis \'ecris le plus grand des deux
nombres (\eng{the bigger number}). La premi\`ere ligne est faite pour toi.}

\vspace{4pt}
\begin{center}
\begin{tabular}{|C{3.4cm}|C{4.6cm}|C{3.4cm}|}
  \hline
  \textbf{The two numbers} & \textbf{Fill with zeros} & \textbf{Bigger} \\ \hline
  $3.3$ and $3.25$ \tall   & $3.3${\color{cborange}\textbf{0}} and $3.25$ & $3.3$ \\ \hline
  $0.6$ and $0.52$ \tall   & $0.6$\zb\ and $0.52$ & \\ \hline
  $4.1$ and $4.09$ \tall   & $4.1$\zb\ and $4.09$ & \\ \hline
  $2.75$ and $2.8$ \tall   & $2.75$ and $2.8$\zb & \\ \hline
  $6.07$ and $6.7$ \tall   & $6.07$ and $6.7$\zb & \\ \hline
  $0.5$ and $0.499$ \tall  & $0.5$\zb\zb\ and $0.499$ & \\ \hline
\end{tabular}
\end{center}

% -- X3 ---------------------------------------------------------------------
\qhead[cbgreen]{X3}{Put them in order}

\noindent \textbf{Step 1.} Fill with zeros, so every number has three digits after the
point.
\fr{\'Etape 1. Ajoute des z\'eros pour que chaque nombre ait trois chiffres apr\`es
la virgule.}

\vspace{8pt}
\noindent\hspace*{1.6em}$2.4$\zb\zb \qquad\quad $2.38$\zb \qquad\quad $2.405$
\qquad\quad $2.04$\zb

\vspace{10pt}
\noindent \textbf{Step 2.} Write the four numbers \textbf{in order of size, starting
with the smallest}.
\fr{\'Etape 2. Range les quatre nombres dans l'ordre croissant, du plus petit au
plus grand (\eng{starting with the smallest}).}

\vspace{8pt}
\noindent\hspace*{1.6em}\blank[2.6cm] \quad \blank[2.6cm] \quad \blank[2.6cm] \quad
\blank[2.6cm]

% -- X4 ---------------------------------------------------------------------
\qhead[cbgreen]{X4}{The snail race}

\noindent Three snails race 1 metre across Mr Bowen's desk. These are their times.
\fr{Trois escargots font la course sur un m\`etre, sur le bureau de Mr Bowen. Voici
leurs temps.}

\begin{wordhelp}[Word help]
  \eng{fastest} = the \textbf{smallest} time. \quad \eng{slowest} = the
  \textbf{biggest} time.
  \fr{\eng{fastest} = le plus rapide = le temps le plus court. \quad \eng{slowest} =
  le plus lent = le temps le plus long.}
\end{wordhelp}

\vspace{4pt}
\begin{center}
\begin{tabular}{|C{2.4cm}|C{3.0cm}|C{4.0cm}|}
  \hline
  \textbf{Snail} & \textbf{Time (min)} & \textbf{Fill with zeros} \\ \hline
  A \tall & 8.5  & $8.5$\zb \\ \hline
  B \tall & 8.25 & $8.25$ \\ \hline
  C \tall & 8.05 & $8.05$ \\ \hline
\end{tabular}
\end{center}

\vspace{4pt}
\begin{enumerate}[label=(\alph*), itemsep=8pt]
  \item Which snail was the \textbf{fastest}? \frin{le plus rapide} \hfill
        \blank[2.6cm]
  \item Which snail was the \textbf{slowest}? \frin{le plus lent} \hfill
        \blank[2.6cm]
\end{enumerate}

% ===========================================================================
% PART 2 -- LINE UP THE POINTS
% ===========================================================================
\hwsection[cbgreen]{Part 2 \textperiodcentered\ Line up the points\frin{Aligne les
virgules}}

\begin{remember}[Adding and subtracting]
  Write the numbers with the \key{points in a line}. Fill the gaps with \key{zeros}.
  Work from the right, like whole numbers. The point in the answer goes
  \key{straight down}.
  \fr{Pose le calcul virgule sous virgule et compl\`ete les cases vides avec des
  z\'eros. Calcule de droite \`a gauche. Dans le r\'esultat, la virgule descend tout
  droit (\eng{straight down}).}
\end{remember}

% -- X5 ---------------------------------------------------------------------
\qhead[cbgreen]{X5}{The point is already there}

\noindent Write a 0 in every orange box. Then add or subtract. The first one is done
for you.
\fr{\'Ecris un 0 dans chaque case orange, puis calcule. Le (a) est fait pour toi.}

\vspace{8pt}
\noindent
\begin{minipage}[t]{0.32\textwidth}\centering
  \textbf{(a)} $4.6 + 2.35$\par\vspace{6pt}
  \colgrid{+}{,4,6,Z}{,2,3,5}{,6,9,5}
\end{minipage}\hfill
\begin{minipage}[t]{0.32\textwidth}\centering
  \textbf{(b)} $12.4 + 3.08$\par\vspace{6pt}
  \colgrid{+}{1,2,4,z}{,3,0,8}{,,,}
\end{minipage}\hfill
\begin{minipage}[t]{0.32\textwidth}\centering
  \textbf{(c)} $7 + 0.45$\par\vspace{6pt}
  \colgrid{+}{,7,z,z}{,0,4,5}{,,,}
\end{minipage}

\vspace{14pt}
\noindent
\begin{minipage}[t]{0.32\textwidth}\centering
  \textbf{(d)} $8.5 - 2.36$\par\vspace{6pt}
  \colgrid{-}{,8,5,z}{,2,3,6}{,,,}
\end{minipage}\hfill
\begin{minipage}[t]{0.32\textwidth}\centering
  \textbf{(e)} $10 - 3.4$\par\vspace{6pt}
  \colgrid{-}{1,0,z,}{,3,4,}{,,,}
\end{minipage}\hfill
\begin{minipage}[t]{0.32\textwidth}\centering
  \textbf{(f)} $1 - 0.36$\par\vspace{6pt}
  \colgrid{-}{,1,z,z}{,0,3,6}{,,,}
\end{minipage}

\vspace{8pt}
\noindent \textbf{(g)} In \textbf{(c)}, where is the point in the number 7?
\frin{O\`u est la virgule dans le nombre 7~?}
\ansline[8.0cm]

% -- X6 ---------------------------------------------------------------------
\qhead[cbgreen]{X6}{Change}

\begin{wordhelp}[Word help]
  \eng{change}: in a shop, the money you get \textbf{back} when you pay too much.
  \; 10 dollars is written \textbf{10.00}.
  \fr{\eng{change} = la monnaie~: l'argent qu'on te rend quand tu donnes plus que le
  prix. 10 dollars s'\'ecrit 10.00.}
\end{wordhelp}

\vspace{4pt}
\noindent
\begin{minipage}[t]{0.48\textwidth}
  \textbf{(a)} Mr Bowen pays for a 6.45-dollar book with a 10-dollar note. How much
  change does he get?
  \fr{Mr Bowen paie un livre \`a 6.45 dollars avec un billet de 10 dollars. Combien
  de monnaie (\eng{change}) lui rend-on~?}
  \par\vspace{8pt}\centering
  \colgrid{-}{1,0,z,z}{,6,4,5}{,,,}
\end{minipage}\hfill
\begin{minipage}[t]{0.48\textwidth}
  \textbf{(b)} Mr Bowen pays for 12.80-dollar shoes with a 20-dollar note. How much
  change does he get?
  \fr{Mr Bowen paie des chaussures \`a 12.80 dollars avec un billet de 20 dollars.
  Combien de monnaie lui rend-on~?}
  \par\vspace{8pt}\centering
  \colgrid{-}{2,0,z,z}{1,2,8,0}{,,,}
\end{minipage}

% -- X7 ---------------------------------------------------------------------
% The scaffold comes off: the grid is still there, the digits are not. Where
% each digit goes is now the student's job, and the point column is the only
% guide.
\qhead[cbgreen]{X7}{Now you line them up}

% A minipage, so the instruction cannot be left at the foot of a page with the
% grids overleaf (it was, on the first build).
\noindent\begin{minipage}{\textwidth}
Write the numbers in the grid yourself. The point goes in the point column. Write a
0 in any gap after the point. Then work it out.
\fr{\'Ecris toi-m\^eme les nombres dans la grille, la virgule dans sa colonne.
Ajoute un 0 dans chaque case vide apr\`es la virgule, puis calcule.}

\vspace{8pt}
\noindent
\begin{minipage}[t]{0.32\textwidth}\centering
  \textbf{(a)} $5.07 + 12.6$\par\vspace{6pt}
  \colgrid{+}{,,,}{,,,}{,,,}
\end{minipage}\hfill
\begin{minipage}[t]{0.32\textwidth}\centering
  \textbf{(b)} $9 - 4.25$\par\vspace{6pt}
  \colgrid{-}{,,,}{,,,}{,,,}
\end{minipage}\hfill
\begin{minipage}[t]{0.32\textwidth}\centering
  \textbf{(c)} $0.8 + 3.45$\par\vspace{6pt}
  \colgrid{+}{,,,}{,,,}{,,,}
\end{minipage}
\end{minipage}

% ===========================================================================
% PART 3 -- POINT OUT, POINT BACK
% ===========================================================================
\hwsection[cbgreen]{Part 3 \textperiodcentered\ Point out, point back\frin{Enl\`eve
la virgule, puis remets-la}}

\begin{remember}[Multiplying a decimal]
  $3.2 \times 7$. \; Take the point out: $32 \times 7 = 224$. \; There is
  \key{1 decimal place} in $3.2$, so put \key{1} back: $22.4$. \\
  \key{Estimate} to check: $3 \times 7 = 21$. \; $22.4$ is close to 21. \checkmark
  \fr{Multiplie sans la virgule, puis remets-la~: le r\'esultat a autant de chiffres
  apr\`es la virgule (\eng{decimal places}) que le nombre de d\'epart. V\'erifie
  avec une estimation (\eng{estimate}).}
\end{remember}

% -- X8 ---------------------------------------------------------------------
% The scaffold fades: row 1 done, rows 2-3 with the calculations printed and
% the numbers missing, rows 4-6 empty.
\qhead[cbgreen]{X8}{Four steps}

\noindent Fill in the table. The first one is done for you.
\fr{Compl\`ete le tableau. La premi\`ere ligne est faite pour toi.}

\vspace{4pt}
\noindent
\begin{tabularx}{\textwidth}{|C{2.3cm}|L|C{2.3cm}|C{2.1cm}|L|}
  \hline
  \textbf{Calculation} & \textbf{Point out} & \textbf{Decimal places} &
  \textbf{Point back} & \textbf{Estimate} \\
  \hline
  $2.4 \times 3$ \tall & $24 \times 3 = 72$ & 1 & $7.2$ & $2 \times 3 = 6$ \\ \hline
  $1.6 \times 4$ \tall & $16 \times 4 =$ & & & $2 \times 4 =$ \\ \hline
  $0.9 \times 6$ \tall & $9 \times 6 =$ & & & $1 \times 6 =$ \\ \hline
  $3.25 \times 2$ \tall & & & & \\ \hline
  $5.3 \times 7$ \tall & & & & \\ \hline
  $4.05 \times 3$ \tall & & & & \\ \hline
\end{tabularx}

% -- X9 ---------------------------------------------------------------------
\qhead[cbgreen]{X9}{The estimate decides}

\noindent All three answers have the same digits. Work out the estimate, then
\textbf{circle} the answer that is close to it.
\fr{Les trois r\'eponses ont les m\^emes chiffres. Fais une estimation, puis entoure
(\eng{circle}) la r\'eponse la plus proche.}

\vspace{4pt}
\noindent
\begin{tabularx}{\textwidth}{|C{2.6cm}|L|C{1.8cm}C{1.8cm}C{1.8cm}|}
  \hline
  \textbf{Calculation} & \textbf{Estimate} &
  \multicolumn{3}{c|}{\textbf{Circle one}} \\
  \hline
  $6.8 \times 4$ \tall  & $7 \times 4 =$ & $2.72$ & $27.2$ & $272$ \\ \hline
  $3.9 \times 6$ \tall  & & $2.34$ & $23.4$ & $234$ \\ \hline
  $12.5 \times 3$ \tall & & $3.75$ & $37.5$ & $375$ \\ \hline
  $0.8 \times 9$ \tall  & & $0.72$ & $7.2$ & $72$ \\ \hline
\end{tabularx}

% -- X10 --------------------------------------------------------------------
% X8's steps with a science label. Every mass is whole and every strength of
% gravity is a decimal from the 3.1 deck, so it is never decimal x decimal.
\qhead[cbgreen]{X10}{Lighter on the Moon}

\noindent On the Moon, weight $=$ mass $\times$ \textbf{1.6}. On Mars, weight $=$ mass
$\times$ \textbf{3.7}. Fill in the steps.
\fr{Sur la Lune, poids (\eng{weight}) $=$ masse $\times$ 1.6~; sur Mars, poids $=$
masse $\times$ 3.7. Compl\`ete les \'etapes.}

\vspace{4pt}
\begin{enumerate}[label=(\alph*), itemsep=10pt]
  \item A 5 kg bag of rice on the \textbf{Moon}: \; $5 \times 1.6$
        \par\vspace{4pt}Point out: \; $5 \times 16 =$ \blank[1.8cm] \qquad
        Point back: \; \blank[1.8cm] N
  \item A 3 kg cat on the \textbf{Moon}: \; $3 \times 1.6$
        \par\vspace{4pt}Point out: \; $3 \times 16 =$ \blank[1.8cm] \qquad
        Point back: \; \blank[1.8cm] N
  \item A 7 kg bag on \textbf{Mars}: \; $7 \times 3.7$
        \par\vspace{4pt}Point out: \; \blank[3.0cm] \qquad
        Point back: \; \blank[1.8cm] N
\end{enumerate}

% ===========================================================================
% PART 4 -- THE POINT GOES STRAIGHT UP
% ===========================================================================
\hwsection[cbgreen]{Part 4 \textperiodcentered\ The point goes straight up\frin{La
virgule monte tout droit}}

\begin{remember}[Dividing a decimal]
  Use short division. The point in the answer goes \key{straight up}. \\
  If a digit is too small to divide, write \key{0} in the answer and carry. \\
  Check with the \key{inverse calculation}: \; $8.24 \div 4 = 2.06$, because $2.06
  \times 4 = 8.24$.
  \fr{Pose la division~: dans le quotient, la virgule monte tout droit
  (\eng{straight up}). Si un chiffre est trop petit pour \^etre divis\'e, \'ecris 0.
  V\'erifie avec l'op\'eration inverse (\eng{inverse calculation}).}
\end{remember}

% -- X11 ---------------------------------------------------------------------
% (c) and (d) need a 0 in the answer; (f) needs a 0 in the orange box under
% the line first, then keeps going.
\qhead[cbgreen]{X11}{One digit in each box}

\noindent The point is already in the answer. Write one digit in each box. In
\textbf{(f)}, write a 0 in the orange box first. Then check.
\fr{La virgule est d\'ej\`a plac\'ee dans le r\'esultat. \'Ecris un chiffre par
case. Dans \textbf{(f)}, \'ecris d'abord un 0 dans la case orange. Ensuite,
v\'erifie.}

\vspace{8pt}
\noindent
\begin{minipage}[t]{0.32\textwidth}\centering
  \textbf{(a)} $6.84 \div 2$\par\vspace{6pt}
  \bus{2}{6,.,8,4}\par\vspace{6pt}
  \small Check: \blank[1.3cm] $\times\, 2 = 6.84$
\end{minipage}\hfill
\begin{minipage}[t]{0.32\textwidth}\centering
  \textbf{(b)} $9.45 \div 3$\par\vspace{6pt}
  \bus{3}{9,.,4,5}\par\vspace{6pt}
  \small Check: \blank[1.3cm] $\times\, 3 = 9.45$
\end{minipage}\hfill
\begin{minipage}[t]{0.32\textwidth}\centering
  \textbf{(c)} $4.28 \div 4$\par\vspace{6pt}
  \bus{4}{4,.,2,8}\par\vspace{6pt}
  \small Check: \blank[1.3cm] $\times\, 4 = 4.28$
\end{minipage}

\vspace{14pt}
\noindent
\begin{minipage}[t]{0.32\textwidth}\centering
  \textbf{(d)} $7.2 \div 8$\par\vspace{6pt}
  \bus{8}{7,.,2}\par\vspace{6pt}
  \small Check: \blank[1.3cm] $\times\, 8 = 7.2$
\end{minipage}\hfill
\begin{minipage}[t]{0.32\textwidth}\centering
  \textbf{(e)} $5.6 \div 7$\par\vspace{6pt}
  \bus{7}{5,.,6}\par\vspace{6pt}
  \small Check: \blank[1.3cm] $\times\, 7 = 5.6$
\end{minipage}\hfill
\begin{minipage}[t]{0.32\textwidth}\centering
  \textbf{(f)} $2.6 \div 4$\par\vspace{6pt}
  \bus{4}{2,.,6,z}\par\vspace{6pt}
  \small Check: \blank[1.3cm] $\times\, 4 = 2.6$
\end{minipage}

% -- X12 --------------------------------------------------------------------
% The point is no longer printed in the answer row. (b) needs the orange 0
% under the line; (c) needs two zeros in the answer: 0.08.
\qhead[cbgreen]{X12}{Now you put the point up}

\noindent\begin{minipage}{\textwidth}
This time the point is \textbf{not} in the answer. Put it there yourself --- straight
up. Then fill the boxes.
\fr{Cette fois, la virgule n'est \textbf{pas} dans le r\'esultat~: place-la
toi-m\^eme, tout droit au-dessus de celle du nombre. Puis remplis les cases.}

\vspace{8pt}
\noindent
\begin{minipage}[t]{0.32\textwidth}\centering
  \textbf{(a)} $8.16 \div 4$\par\vspace{6pt}
  \bus[np]{4}{8,.,1,6}
\end{minipage}\hfill
\begin{minipage}[t]{0.32\textwidth}\centering
  \textbf{(b)} $9.3 \div 6$\par\vspace{6pt}
  \bus[np]{6}{9,.,3,z}
\end{minipage}\hfill
\begin{minipage}[t]{0.32\textwidth}\centering
  \textbf{(c)} $0.72 \div 9$\par\vspace{6pt}
  \bus[np]{9}{0,.,7,2}
\end{minipage}
\end{minipage}

% -- X13 --------------------------------------------------------------------
% The mean is an add and a divide, and the divide is X11. Both data sets are
% from Science 3.1: the falling-ball lab and the ruler catch.
\qhead[cbgreen]{X13}{The mean, in two steps}

\begin{remember}[The mean]
  \key{Step 1:} add up all the values. \quad \key{Step 2:} divide by how many there
  are.
  \fr{La moyenne (\eng{mean})~: additionne toutes les valeurs, puis divise par le nombre de valeurs.}
\end{remember}

\vspace{2pt}
\begin{enumerate}[label=(\alph*), itemsep=10pt]
  \item Mr Bowen drops a ball three times and times it. The times are
        \textbf{0.46 s}, \textbf{0.44 s} and \textbf{0.45 s}. Find the mean.
        \fr{Mr Bowen l\^ache une balle trois fois et chronom\`etre chaque chute.
        Calcule la moyenne.}
        \par\vspace{6pt}\hspace*{1.6em}Step 1: \; $0.46 + 0.44 + 0.45 =$ \blank[2.0cm]
        \par\vspace{6pt}\hspace*{1.6em}Step 2: \; \blank[2.0cm] $\div\ 3 =$
        \blank[2.0cm] s
  \item His reaction times in the ruler catch are \textbf{0.21 s}, \textbf{0.19 s},
        \textbf{0.23 s} and \textbf{0.17 s}. Find the mean.
        \fr{Voici ses temps de r\'eaction (\eng{reaction times}) au test de la
        r\`egle. Calcule la moyenne.}
        \par\vspace{6pt}\hspace*{1.6em}Step 1: \; \blank[6.0cm] $=$ \blank[2.0cm]
        \par\vspace{6pt}\hspace*{1.6em}Step 2: \; \blank[2.0cm] $\div$ \nbox $=$
        \blank[2.0cm] s
\end{enumerate}

% ===========================================================================
% PART 5 -- MAKING IT EASIER
% ===========================================================================
\hwsection[cbgreen]{Part 5 \textperiodcentered\ Making it easier\frin{Calculer plus facilement}}

% -- X14 --------------------------------------------------------------------
\qhead[cbgreen]{X14}{``Of'' means $\times$}

\begin{remember}[Half, a quarter, a tenth]
  \key{half of} = $\div\, 2$ = $\times\, 0.5$ \qquad
  \key{a quarter of} = $\div\, 4$ = $\times\, 0.25$ \qquad
  \key{a tenth of} = $\div\, 10$ = $\times\, 0.1$
  \fr{\eng{half of} = la moiti\'e de \quad \eng{a quarter of} = le quart de \quad \eng{a tenth of} = le dixi\`eme de}
\end{remember}

\vspace{2pt}
\noindent Fill in the table. The first one is done for you.
\fr{Compl\`ete le tableau. La premi\`ere ligne est faite pour toi.}

\vspace{4pt}
\begin{center}
\begin{tabular}{|C{4.0cm}|C{3.6cm}|C{2.6cm}|}
  \hline
  \textbf{In English} & \textbf{Calculation} & \textbf{Answer} \\ \hline
  half of 18 \tall       & $18 \div 2$ & $9$ \\ \hline
  a quarter of 20 \tall  & $20 \div$ \nbox & \\ \hline
  a tenth of 35 \tall    & $35 \div$ \nbox & \\ \hline
  half of 7 \tall        & & \\ \hline
  a quarter of 6 \tall   & & \\ \hline
\end{tabular}
\end{center}

% -- X15 --------------------------------------------------------------------
\qhead[cbgreen]{X15}{Times 9 and times 11}

\begin{remember}[Start from $\times\, 10$]
  $\times\, 9$: do $\times\, 10$, then \key{take away} one lot. \qquad
  $\times\, 11$: do $\times\, 10$, then \key{add} one lot.
  \fr{$\times\, 9$~: multiplie par 10, puis retire le nombre une fois. \quad
  $\times\, 11$~: multiplie par 10, puis ajoute-le une fois.}
\end{remember}

\vspace{4pt}
\begin{enumerate}[label=(\alph*), itemsep=12pt]
  \item $2.6 \times 9 = 26 - 2.6 = \blank[2.2cm]$
  \item $1.4 \times 9 = \blank[1.6cm] - \blank[1.6cm] = \blank[2.2cm]$
  \item $2.6 \times 11 = 26 + 2.6 = \blank[2.2cm]$
  \item $3.2 \times 11 = \blank[1.6cm] + \blank[1.6cm] = \blank[2.2cm]$
\end{enumerate}

% -- X16 --------------------------------------------------------------------
\qhead[cbgreen]{X16}{Split the number}

\noindent \textbf{Partitioning} means splitting a number into easy parts:
$4.3 = 4 + 0.3$.
\fr{\eng{Partitioning}~: d\'ecomposer un nombre en parties faciles \`a calculer.}

\vspace{6pt}
\begin{enumerate}[label=(\alph*), itemsep=12pt]
  \item $4.2 \times 3 = 4 \times 3 + 0.2 \times 3 = \blank[1.6cm] + \blank[1.6cm] =
        \blank[2.0cm]$
  \item $5.3 \times 4 = 5 \times 4 + 0.3 \times 4 = \blank[1.6cm] + \blank[1.6cm] =
        \blank[2.0cm]$
  \item $7.4 \times 2 = \blank[1.6cm] + \blank[1.6cm] = \blank[2.0cm]$
\end{enumerate}

% ===========================================================================
% PART 6 -- THE WORDS THAT CHOOSE
% ===========================================================================
\hwsection[cbgreen]{Part 6 \textperiodcentered\ The words that choose\frin{Les mots
qui indiquent l'op\'eration}}

% -- X17 --------------------------------------------------------------------
\qhead[cbgreen]{X17}{Which operation?}

\noindent Write $+$, $-$, $\times$ or $\div$ for each one. The first one is done for
you.
\fr{Pour chaque expression, \'ecris $+$, $-$, $\times$ ou $\div$. La premi\`ere
ligne est faite pour toi.}

\vspace{4pt}
% The words column is wider in the French copy, which adds a gloss to each row.
\newlength{\wordcol}
\iffr\setlength{\wordcol}{9.4cm}\else\setlength{\wordcol}{6.2cm}\fi
\begin{center}
\begin{tabular}{|p{\wordcol}|C{2.2cm}|}
  \hline
  \textbf{The words} & \textbf{Operation} \\ \hline
  altogether \frin{en tout} \tall & $+$ \\ \hline
  change \frin{la monnaie} \tall & \\ \hline
  share equally \frin{partager \'equitablement} \tall & \\ \hline
  how much do 4 cost? \frin{combien co\^utent 4 objets~?} \tall & \\ \hline
  how much heavier? \frin{de combien plus lourd~?} \tall & \\ \hline
  the total \frin{le total} \tall & \\ \hline
\end{tabular}
\end{center}

% -- X18 --------------------------------------------------------------------
\qhead[cbgreen]{X18}{Calculation first}

\noindent Write the calculation in the boxes: number, sign, number. Then the answer.
\fr{\'Ecris le calcul~: un nombre, un signe dans la case, un nombre. Puis le
r\'esultat.}

\begin{enumerate}[label=(\alph*), itemsep=8pt]
  \item A bottle of water costs 0.85 dollars. Mr Bowen buys 4. How much does he pay?
        \fr{Une bouteille d'eau co\^ute 0.85 dollar. Mr Bowen en ach\`ete 4. Combien paie-t-il~?}
        \par\vspace{6pt}\hspace*{1.6em}\blank[2.0cm] \nbox \blank[1.6cm] $=$
        \blank[2.2cm] dollars
  \item Mr Bowen pays for a 2.65-dollar juice with a 5-dollar note. How much change
        does he get?
        \fr{Mr Bowen paie un jus de fruits \`a 2.65 dollars avec un billet de 5
        dollars. Combien de monnaie (\eng{change}) lui rend-on~?}
        \par\vspace{6pt}\hspace*{1.6em}\blank[2.0cm] \nbox \blank[1.6cm] $=$
        \blank[2.2cm] dollars
  \item Mr Bowen and three friends share 9.20 dollars equally. How much does each
        person get?
        \fr{Mr Bowen et trois amis se partagent \'equitablement (\eng{share
        equally}) 9.20 dollars. Combien chacun re\c{c}oit-il~?}
        \par\vspace{6pt}\hspace*{1.6em}How many people? \nbox \hfill
        \blank[2.0cm] \nbox \blank[1.6cm] $=$ \blank[2.2cm] dollars
  \item A bag of sugar has a mass of 0.5 kg. What is its \textbf{weight} on Earth?
        Weight $=$ mass $\times$ 10.
        \fr{Un sac de sucre a une masse de 0.5 kg. Quel est son poids (\eng{weight}) sur Terre~?}
        \par\vspace{6pt}\hspace*{1.6em}\blank[2.0cm] \nbox \blank[1.6cm] $=$
        \blank[2.2cm] N
\end{enumerate}

% -- X19 --------------------------------------------------------------------
% The sheet in miniature: a multiplication with the point taken out and put
% back, then a subtraction with the zeros written. Last on purpose.
% In the French copy it starts the last page, so its working box is not left
% on a page of its own (it was).
\iffr\clearpage\fi
\qhead[cbgreen]{X19}{The whole sheet in one question}

\noindent Mr Bowen buys 3 mangoes at 1.45 dollars each. He pays with a 5-dollar note.
How much change does he get?
\fr{Mr Bowen ach\`ete 3 mangues \`a 1.45 dollar pi\`ece. Il paie avec un billet de 5
dollars. Combien de monnaie lui rend-on~?}

\vspace{8pt}
\noindent\hspace*{1.6em}\textbf{Step 1.} The mangoes: \quad $145 \times 3 =$
\blank[1.8cm], \; so $1.45 \times 3 =$ \blank[1.8cm] dollars.
\fr{\'Etape 1~: le prix des mangues.}

\vspace{10pt}
\noindent\hspace*{1.6em}\textbf{Step 2.} The change: \quad $5.00 \;-$ \blank[1.8cm]
$=$ \blank[1.8cm] dollars.
\fr{\'Etape 2~: la monnaie.}

\vspace{12pt}
\noindent Work it out here. Write $5$ as $5.00$ and line up the points.
\fr{Calcule ici. \'Ecris 5 sous la forme 5.00 et pose le calcul virgule sous
virgule.}

% Taller in the English copy, which has the room on its last page.
\newlength{\finalboxht}
\iffr\setlength{\finalboxht}{8.0cm}\else\setlength{\finalboxht}{7.0cm}\fi
\begin{workspace}[\finalboxht]\end{workspace}

% ===========================================================================
% ANSWER KEY -- TEACHER COPY ONLY
% ===========================================================================
\ifkey
\clearpage

\hwsection[cbred]{Answer Key --- Teacher Copy}

\linespread{1.0}\selectfont
\setlist[enumerate]{itemsep=2pt, topsep=3pt, leftmargin=*}
\setlist[itemize]{itemsep=2pt, topsep=3pt, leftmargin=*}

\qhead[cbred]{X1}{Two parts of a number}
\noindent $14.6$: 14 and .6 \sep $0.38$: 0 and .38 \sep $7.05$: 7 and .05 \sep
$120.9$: 120 and .9

\smallskip
\noindent{\small $0.38$ is the row that gets a blank: there is a whole-number part, and
it is 0. $7.05$ is the row that gets ``.5''. Both are the same habit --- not writing a
zero that is holding a place --- and X2 depends on it.}

\qhead[cbred]{X2}{Fill with zeros, then compare}
\noindent $0.60$, bigger $0.6$ \sep $4.10$, bigger $4.1$ \sep $2.80$, bigger $2.8$ \sep
$6.70$, bigger $6.7$ \sep $0.500$, bigger $0.5$

\smallskip
\noindent{\small \textbf{In every row the bigger number is the shorter one}, and that
is on purpose: it is the opposite of what ``more digits'' suggests, and the zeros are
what show it. \\
$6.07$ against $6.7$ is the row to discuss: both have a 0 and a 7, and it is
\emph{where} the 0 is that decides. The 0 in $6.07$ holds the tenths open; the 0 in
$6.70$ adds nothing. \\
$0.5$ against $0.499$ needs two boxes. A student who writes $0.50$ and compares $50$
with $499$ has still been fooled by length.}

\qhead[cbred]{X3}{Put them in order}
\noindent Step 1: $2.400, \; 2.380, \; 2.405, \; 2.040$ \\
Step 2: $2.04, \; 2.38, \; 2.4, \; 2.405$

\smallskip
\noindent{\small Accept the numbers written with or without their filled-in zeros. \\
$2.4$ against $2.405$ is the close pair: $2.400 < 2.405$, so $2.4$ comes first. Expect
$2.405$ before $2.4$ from a student who stopped filling at two places. $2.04$ first is
the one most will get, because the zero is so visible.}

\qhead[cbred]{X4}{The snail race}
\noindent $8.50$ \sep (a) \textbf{C} ($8.05$) \sep (b) \textbf{A} ($8.5$)

\smallskip
\noindent{\small The slowest snail has the shortest-looking time. Expect B for (b) from
anyone comparing length, and A for (a) from anyone who has read ``fastest'' as
``biggest''. The word help is on the page; send them back to it.}

\qhead[cbred]{X5}{The point is already there}
\noindent (b) $12.40 + 3.08 = \mathbf{15.48}$ \quad (c) $7.00 + 0.45 = \mathbf{7.45}$
\quad (d) $8.50 - 2.36 = \mathbf{6.14}$ \\
(e) $10.0 - 3.4 = \mathbf{6.6}$ \quad (f) $1.00 - 0.36 = \mathbf{0.64}$ \\
(g) After the 7: \; $7 = 7.00$. \; (Accept ``at the end'', ``after the 7''.)

\smallskip
\noindent{\small The grid makes the alignment impossible to get wrong, so what is left
is the subtraction with zeros. (d) and (f) need a borrow from a 0 the student wrote
themselves; expect $6.26$ for (d) and $1.36$ for (f) --- ``take the smaller digit from
the bigger'' in each column, which a 0 on top invites. \\
(g) is the question the grid answers by itself. A whole number has its point at the
end, and that is why (c)'s 7 is in the ones column and not under the 4.}

\qhead[cbred]{X6}{Change}
\noindent (a) $10.00 - 6.45 = \mathbf{3.55}$ dollars \quad (b) $20.00 - 12.80 =
\mathbf{7.20}$ dollars

\smallskip
\noindent{\small (b) answers in money, so $7.20$, not $7.2$ (HW 9's E3). \\
(a) borrows across three zeros and is the hardest subtraction on the sheet. If it goes
wrong, check it column by column from the right; the inverse ($3.55 + 6.45 = 10.00$)
is the fastest way for the student to see it.}

\qhead[cbred]{X7}{Now you line them up}
\noindent (a) $5.07 + 12.60 = \mathbf{17.67}$ \quad (b) $9.00 - 4.25 = \mathbf{4.75}$
\quad (c) $0.80 + 3.45 = \mathbf{4.25}$

\smallskip
\noindent{\small \textbf{Look at where the digits went before marking the answer.}
(a) is the test: $12.6$ needs the tens column and $5.07$ does not, so a student who
fills the grid from the left puts the 5 under the 1 and gets $62.67$ or similar. \\
(b): the 9 goes in the ones, with two zeros after the point. Expect $9 - 4.25$
written as $9.25 - 4$ or with the 9 in the hundredths.}

\qhead[cbred]{X8}{Four steps}
\noindent
\begin{tabular}{@{}lllll@{}}
  $1.6 \times 4$ & $16 \times 4 = 64$ & 1 & $6.4$ & $2 \times 4 = 8$ \\
  $0.9 \times 6$ & $9 \times 6 = 54$ & 1 & $5.4$ & $1 \times 6 = 6$ \\
  $3.25 \times 2$ & $325 \times 2 = 650$ & 2 & $6.50 = 6.5$ & $3 \times 2 = 6$ \\
  $5.3 \times 7$ & $53 \times 7 = 371$ & 1 & $37.1$ & $5 \times 7 = 35$ \\
  $4.05 \times 3$ & $405 \times 3 = 1215$ & 2 & $12.15$ & $4 \times 3 = 12$ \\
\end{tabular}

\smallskip
\noindent{\small The third column is the step, and it counts decimal places in the
\emph{question}, not digits. Expect 3 for $3.25$ from anyone counting all the digits;
the answer then comes out $0.650$ and the estimate of 6 throws it out. \\
$4.05$: the 0 is a digit and stays in when the point comes out --- $405$, not $45$. Expect
$45 \times 3 = 135$, $1.35$, from anyone who drops it. The estimate (12) catches that too.}

\qhead[cbred]{X9}{The estimate decides}
\noindent $27.2$ ($7 \times 4 = 28$) \sep $23.4$ ($4 \times 6 = 24$) \sep $37.5$
($12 \times 3 = 36$, or $13 \times 3 = 39$) \sep $7.2$ ($1 \times 9 = 9$)

\smallskip
\noindent{\small No multiplication is needed in this question and that is the point
of it. A student who works out $68 \times 4$ first has done more than was asked and
should be told so --- warmly. \\
$0.8 \times 9$: the estimate $1 \times 9 = 9$ is a bit big, and $7.2$ is the only
choice near it. Some will circle $0.72$ because ``$0.8$ is less than 1''; ask whether
nine lots of nearly 1 could be less than 1.}

\qhead[cbred]{X10}{Lighter on the Moon}
\noindent (a) $80$, so $\mathbf{8}$ N ($8.0$) \quad (b) $48$, so $\mathbf{4.8}$ N \quad
(c) $7 \times 37 = 259$, so $\mathbf{25.9}$ N

\smallskip
\noindent{\small (a) comes back as $8.0$, which is $8$; accept both (nothing was
rounded). Expect $80$ from a student who skipped putting the point back --- a bag of
rice weighing 80 N on the Moon would weigh more than on Earth (50 N), which is worth
asking about. \\
This is HW 9's F2 with smaller numbers and the steps printed.}

\qhead[cbred]{X11}{One digit in each box}
\noindent (a) $3.42$ \sep (b) $3.15$ \sep (c) $\mathbf{1.07}$ \sep (d) $\mathbf{0.9}$
\sep (e) $\mathbf{0.8}$ \sep (f) $\mathbf{0.65}$ \\
Checks: $3.42 \times 2$, $3.15 \times 3$, $1.07 \times 4$, $0.9 \times 8$, $0.8 \times
7$, $0.65 \times 4$.

\smallskip
\noindent{\small (c), (d) and (e) each need a 0 written in an answer box. The frame
has a box there, so a blank box is visible --- that is the scaffold doing its job.
Expect $1.7$ for (c) if the student skipped the box. \\
(d) and (e) start with ``8 into 7 will not go'': the 0 goes in the ones, then the
remainder carries. \\
(f) runs out of digits: $2.60$, and the orange 0 is where the last 5 comes from.
Expect ``$0.6$ remainder 2''.}

\qhead[cbred]{X12}{Now you put the point up}
\noindent (a) $\mathbf{2.04}$ \quad (b) $9.30 \div 6 = \mathbf{1.55}$ \quad
(c) $\mathbf{0.08}$

\smallskip
\noindent{\small Look for the point in the answer row first; it should sit exactly
above the dividend's point. \\
(c) is the hardest item in Part 4: $0 \div 9 = 0$, then $7 \div 9$ will not go, so a
second 0, then $72 \div 9 = 8$. Expect $0.8$ and $8$. \\
(b) needs the orange 0 to finish: $9.3 \div 6$ stops at $1.5$ remainder 3, and $30
\div 6 = 5$ gives the last digit.}

\qhead[cbred]{X13}{The mean, in two steps}
\noindent (a) $0.46 + 0.44 + 0.45 = 1.35$, \; $1.35 \div 3 = \mathbf{0.45}$ s \\
(b) $0.21 + 0.19 + 0.23 + 0.17 = 0.80$, \; $0.80 \div 4 = \mathbf{0.2}$ s

\smallskip
\noindent{\small (a) is the drop time from HW 9's F3, so a student doing both sheets
meets $0.45$ s twice --- on purpose. \\
(b): ``how many there are'' is 4, the box in Step 2. Expect $\div\ 2$ from a student who
counted the digits after the point. The answer, $0.2$ s, is the reaction time the
Science 3.1 deck quoted.}

\qhead[cbred]{X14}{``Of'' means $\times$}
\noindent a quarter of 20: $20 \div 4 = 5$ \sep a tenth of 35: $35 \div 10 = 3.5$ \sep
half of 7: $7 \div 2 = 3.5$ \sep a quarter of 6: $6 \div 4 = 1.5$

\smallskip
\noindent{\small The last two rows have no calculation printed; a student who writes
only the answer has skipped the column. \\
``A quarter of 6'' is the hardest: $6 \div 4$ does not divide exactly, so it is short
division with the point --- $6.0 \div 4 = 1.5$ --- which is Part 4 again.}

\qhead[cbred]{X15}{Times 9 and times 11}
\noindent (a) $23.4$ \quad (b) $14 - 1.4 = 12.6$ \quad (c) $28.6$ \quad (d) $32 + 3.2
= 35.2$

\smallskip
\noindent{\small (a) and (b) are subtractions with a hidden zero: $26.0 - 2.6$. Expect
$24.6$ from $26 - 2 + 0.6$: the 0.6 added instead of taken away. Writing $26.0$ in a
column fixes it.}

\qhead[cbred]{X16}{Split the number}
\noindent (a) $12 + 0.6 = 12.6$ \quad (b) $20 + 1.2 = 21.2$ \quad (c) $7 \times 2 + 0.4
\times 2 = 14 + 0.8 = 14.8$

\smallskip
\noindent{\small $0.3 \times 4 = 1.2$ is the one to watch: it crosses 1. Expect
$0.12$, which is a decimal place counted in the wrong place.}

\qhead[cbred]{X17}{Which operation?}
\noindent change $-$ \sep share equally $\div$ \sep how much do 4 cost? $\times$ \sep how
much heavier? $-$ \sep the total $+$

\smallskip
\noindent{\small ``How much do 4 cost?'' is $\times$, and some will write $+$ --- which
is not wrong, just slow (repeated addition). Accept it with a word.}

\qhead[cbred]{X18}{Calculation first}
\noindent (a) $0.85 \times 4 = 3.40$ dollars \sep (b) $5.00 - 2.65 = 2.35$ dollars \\
(c) \textbf{4} people; $9.20 \div 4 = 2.30$ dollars \sep (d) $0.5 \times 10 = 5$ N

\smallskip
\noindent{\small (c) is HW 9's E2(c) with the box that makes the student count: ``Mr
Bowen and three friends'' is four. A 3 in the people box is the English error, and
this sheet catches it before the division starts. \\
(d) is the one science question on the sheet, and the number comes from Science 3.1:
$W = m \times 10$.}

\qhead[cbred]{X19}{The whole sheet in one question}
\noindent Step 1: $145 \times 3 = 435$, so $\mathbf{4.35}$ dollars. \quad Step 2:
$5.00 - 4.35 = \mathbf{0.65}$ dollars.

\smallskip
\noindent{\small Point out and back (Part 3), then zeros written and subtracted (Part
2), with ``change'' choosing the operation (Part 6). Expect $1.35$ from a student who
subtracted $4.35$ from $5.00$ the easy way round in the tenths (X5(d)'s error).}

\fi
\end{document}
